%% supplement.tex — Supplementary Material
%% Myhill-Nerode Quotient Learning via Kernel MMD (continuous spaces)


\documentclass[10pt]{article}

%% --- Page geometry ---
\usepackage[margin=1in]{geometry}

%% --- Core packages ---
\usepackage[utf8]{inputenc}
\usepackage[T1]{fontenc}
\usepackage{amsmath, amssymb, amsthm}
\usepackage{mathtools}
\usepackage{bm}
\usepackage{booktabs}
\usepackage{graphicx}
\usepackage{hyperref}
\usepackage{xcolor}
\usepackage{algorithm}
\usepackage{algpseudocode}
\usepackage{natbib}
\usepackage{url}
\usepackage{enumerate}

%% --- Theorem environments ---
\newtheorem{theorem}{Theorem}[section]
\newtheorem{lemma}[theorem]{Lemma}
\newtheorem{proposition}[theorem]{Proposition}
\newtheorem{corollary}[theorem]{Corollary}
\newtheorem{definition}{Definition}[section]
\newtheorem{remark}{Remark}[section]
\newtheorem{fact}{Fact}[section]

%% --- Math macros (identical to main.tex) ---
\newcommand{\calQ}{\mathcal{Q}}
\newcommand{\calH}{\mathcal{H}}
\newcommand{\calO}{\mathcal{O}}
\newcommand{\calA}{\mathcal{A}}
\newcommand{\calS}{\mathcal{S}}
\newcommand{\calK}{\mathcal{K}}
\newcommand{\calC}{\mathcal{C}}
\newcommand{\calP}{\mathcal{P}}
\newcommand{\calM}{\mathcal{M}}
\newcommand{\calB}{\mathcal{B}}
\newcommand{\bbR}{\mathbb{R}}
\newcommand{\bbE}{\mathbb{E}}
\newcommand{\epsMMD}{\varepsilon_{\mathrm{MMD}}}
\newcommand{\epsmin}{\varepsilon_{\min}}
\newcommand{\Qstar}{\mathcal{Q}^*}
\newcommand{\mmd}{\mathrm{MMD}}
\newcommand{\TV}{\mathrm{TV}}
\newcommand{\RKHS}{\mathcal{H}_k}
\newcommand{\phih}{\phi_h}
\newcommand{\Khist}{K_{\mathrm{hist}}}

%% --- Section labelling as A, B, C, ... ---
\renewcommand{\thesection}{\Alph{section}}
\renewcommand{\thesubsection}{\thesection.\arabic{subsection}}

%% ============================================================
\begin{document}

\begin{center}
  {\LARGE\bfseries Supplementary Material}\\[4pt]
  {\large Myhill--Nerode Quotient Learning\\
          for Continuous Observation--Action Spaces\\
          via Kernel Maximum Mean Discrepancy}\\[4pt]
  {\normalsize Anonymous Author(s)}
\end{center}

\vspace{1em}

This document contains the proofs of Lemma~1, Theorem~1 and Proposition~1 of
the main paper, the exact configuration of every experiment (as run by the
scripts shipped with the code), and per-seed result tables.  Every number is
taken from a result file under \texttt{experiments/*/results/}; the file is
named next to each table.

% ============================================================
\section{Proofs of Lemma~1 and Theorem~1}
\label{sec:proof-main}
% ============================================================

Throughout, $\kappa^2=\sup_x\kappa(x,x)$ and, as required by
\citet[Thm.~7]{Gretton2012}, $0\le\kappa(x,y)\le\kappa^2$; the RBF kernel
satisfies both with $\kappa=1$.  For a distribution $P$ let $\mu_P=\bbE_P[\phi(c)]$
be its mean embedding, so that $\mmd(P,Q)=\|\mu_P-\mu_Q\|_{\RKHS}$.

\subsection{Lemma~1: MMD is dominated by total variation}

\begin{proof}
For $\|f\|_{\RKHS}\le1$, $|f(x)|=|\langle f,\phi(x)\rangle|\le\kappa$, hence
\[
  \mmd(P,Q)=\sup_{\|f\|\le1}\Big|\int f\,d(P-Q)\Big|
  \le\kappa\,\|P-Q\|_1=2\kappa\,\TV(P,Q).
\]
For the second claim, $P_b=\int P^{(k)}_{h',\pi}\,\nu_b(dh')$ where $\nu_b$ is
the conditional law of the history given the bin, so
$\mu_{P_b}=\int\mu_{P_{h'}}\nu_b(dh')$ (Bochner integral; the embedding is
bounded) and
\[
  \|\mu_{P_h}-\mu_{P_b}\|_{\RKHS}\le\int\|\mu_{P_h}-\mu_{P_{h'}}\|\,\nu_b(dh')
  \le2\kappa\int\TV(P_h,P_{h'})\,\nu_b(dh')\le2\kappa L\,w_b .\qedhere
\]
\end{proof}

\subsection{Two facts about the biased MMD statistic}

Let $X=(x_1,\dots,x_m)\sim P^{\otimes m}$ and $Y=(y_1,\dots,y_m)\sim Q^{\otimes m}$
be independent and let $\widehat{\mmd}$ be the square root of the biased
statistic (Eq.~(1) of the main paper, $m_1=m_2=m$).

\begin{fact}[{\citealp[Thm.~7]{Gretton2012}}]
\label{fact:dev}
For every $t>0$,
$\Pr\big[\,|\widehat{\mmd}-\mmd(P,Q)|>4\kappa/\sqrt m+t\,\big]\le2\exp(-t^2m/(4\kappa^2))$.
\end{fact}

\begin{fact}[{\citealp[Cor.~9]{Gretton2012}}]
\label{fact:null}
If $P=Q$ then, for every $\alpha_0\in(0,1)$,
$\Pr\big[\widehat{\mmd}>t_{\alpha_0}\big]\le\alpha_0$ with
$t_{\alpha_0}=\sqrt{2\kappa^2/m}\,\big(1+\sqrt{2\log(1/\alpha_0)}\big)$.
\end{fact}

\subsection{Per-pair power}

\begin{lemma}[Per-pair decision]
\label{lem:pair}
Fix $\alpha_0,\delta'\in(0,1)$ and decide ``different'' iff
$\widehat{\mmd}>t_{\alpha_0}$.  (i) If $P=Q$ the decision is wrong with
probability at most $\alpha_0$.  (ii) If $\mmd(P,Q)\ge\varepsilon>0$ and
\begin{equation}
  m\;\ge\;\frac{16\kappa^2}{\varepsilon^2}\Big(8+\log\frac{2}{\alpha_0\delta'}\Big),
  \label{eq:m-pair}
\end{equation}
the decision is wrong with probability at most $\delta'$.
\end{lemma}

\begin{proof}
(i) is Fact~\ref{fact:null}.  For (ii) put $t=2\kappa\sqrt{\log(2/\delta')/m}$,
so that $2\exp(-t^2m/(4\kappa^2))=\delta'$.  By Fact~\ref{fact:dev}, with
probability at least $1-\delta'$,
$\widehat{\mmd}\ge\varepsilon-4\kappa/\sqrt m-t$, and the decision is correct
as soon as $\varepsilon-4\kappa/\sqrt m-t>t_{\alpha_0}$, i.e.
\[
  \sqrt m\,\varepsilon\;>\;\kappa\Big[4+\sqrt2+2\sqrt{\log(1/\alpha_0)}+2\sqrt{\log(2/\delta')}\Big].
\]
Using $(a+b+c+d)^2\le4(a^2+b^2+c^2+d^2)$ the right-hand side squared is at most
$\kappa^2\big[4(4+\sqrt2)^2+16\log(1/\alpha_0)+16\log(2/\delta')\big]
\le\kappa^2\big[128+16\log(2/(\alpha_0\delta'))\big]$, and \eqref{eq:m-pair}
guarantees $m\varepsilon^2$ exceeds this bound.
\end{proof}

\subsection{Proof of Theorem~1}

\begin{proof}
Let $\calP$ be the set of tested pairs, $|\calP|\le\Khist^2/2$.  Run the
decision of Lemma~\ref{lem:pair} on every pair with
$\alpha_0=\delta'=\delta/\Khist^2$.  With $\varepsilon=\epsMMD$,
\eqref{eq:m-pair} becomes condition~(2) of the main paper,
$m\ge(16\kappa^2/\epsMMD^2)(8+\log(2\Khist^4/\delta^2))$.  By a union bound
the probability that \emph{some} pair is decided wrongly is at most
$|\calP|(\alpha_0+\delta')\le(\Khist^2/2)(2\delta/\Khist^2)=\delta$.  On the
complementary event every same-class pair (in $\Qstar_B$) is merged and no
cross-class pair is merged, so the connected components of the union--find
graph are exactly the classes of $\Qstar_B$, i.e.\ $\hat\calQ_B=\Qstar_B$.

For the trajectory budget, each trajectory contributes $T'=T-\ell-k+1$
$(\text{history},\text{continuation})$ pairs; for $\ell=1$ quantile keys have
equal mass, so each of the $\Khist$ keys receives $NT'/\Khist$ of them (for
$\ell>1$ replace $1/\Khist$ by the smallest key mass).  Under the
$\tau$-mixing assumption of the theorem, thinning the within-trajectory
sequence to windows separated by $\tau+k$ steps leaves
$m\ge NT'/(c\Khist(\tau+k))$ effectively independent samples per key
(consecutive continuation windows overlap over $k$ steps and are decorrelated
after a further $\tau$ steps).  Requiring this to exceed the per-pair bound
gives
\[
  N\;\ge\;\frac{16c\,\kappa^2\,\Khist(\tau+k)}{T'\,\epsMMD^2}
          \Big(8+\log\frac{2\Khist^4}{\delta^2}\Big)
   \;=\;O\!\Big(\frac{\Khist(\tau+k)\log(\Khist/\delta)}{T\,\epsMMD^2}\Big)
\]
for $T\gg\ell+k$.  If the implementation caps each side at $m_{\max}$
samples, the theorem applies with $m=\min(m,m_{\max})$, i.e.\ it requires
$m_{\max}$ to exceed the right-hand side of \eqref{eq:m-pair}.
\end{proof}

\begin{remark}[Why the constants matter]
With $\kappa=1$, $\Khist=8$, $\delta=0.05$ the per-pair requirement is
$m\ge(16/\epsMMD^2)(8+\log(2\cdot8^4/0.05^2))\approx 400/\epsMMD^2$.  At the
sample cap $m_{\max}=300$ used in the experiments this guarantees detection
only for $\epsMMD\gtrsim1.15$; the permutation test is far less conservative
than the analytic threshold in practice, but the calculation shows that
distinctions with small MMD gaps are \emph{not} resolved at $m_{\max}=300$,
which is the mechanism behind Remark~4 of the main paper.
\end{remark}

% ============================================================
\section{Proof of Proposition~1}
\label{sec:proof-prop1}
% ============================================================

Setting: $\ell=1$; a history is keyed by the bin of the scalar $z=z(o_t)$;
the POMDP has finite latent index $n$ with class map $c$; $F$ is the pooled
CDF of $z$ (over trajectories and time steps) and $F_N$ its empirical
counterpart.  Empirical edges are $\hat e_j=F_N^{-1}(j/B)$ and population edges
$e_j=F^{-1}(j/B)$, $j=1,\dots,B-1$.

\paragraph{Glivenko--Cantelli.}
The $N$ trajectories are i.i.d.\ blocks of $T$ observations, so
$F_N(x)=N^{-1}\sum_i G_i(x)$ with $G_i(x)=T^{-1}\sum_t\mathbf 1\{z^{(i)}_t\le x\}$
i.i.d.\ in $i$; the class $\{x\mapsto G_i(x)\}$ is monotone and bounded, so
$\sup_x|F_N(x)-F(x)|\to0$ almost surely.

\paragraph{Under (a) and (b).}
Let $G_i=(\sup I_i,\inf I_{i+1})$ be the $i$-th gap; $F$ is constant on
$\overline{G_i}$ at level $q_i=\sum_{j\le i}p_j$, and by (b) $q_i=j_i/B$ for an
integer $j_i$.  Hence the population edge $e_{j_i}$ lies in $\overline{G_i}$ and
every population bin is contained in a single $I_i$: all population bins are
pure, $P_b=P_{c(b)}$, and $\Qstar_B$ is the class partition (each class has
mass $p_i\ge1/B$ by (b), so each class owns at least one bin), i.e.\
$|\Qstar_B|=n$.

For the empirical edges, $F_N(\sup I_i)-q_i=O_P(N^{-1/2})$ ($T$ fixed), so
the number of observations of class $i$ (resp.\ $i{+}1$) that fall on the wrong
side of $\hat e_{j_i}$ is $O_P(T\sqrt N)$, while each bin holds $NT/B$
observations.  Thus the foreign-class \emph{fraction} of every empirical bin is
$O_P(B/\sqrt N)\to0$ (asymptotic purity), and by the mixture identity
$\mmd(\lambda P+(1-\lambda)P',P)=(1-\lambda)\mmd(P,P')$ the MMD between an
empirical bin law and its class law is $O_P(B/\sqrt N)$.

Now run variant~(T) with level $\alpha_N\to0$, $\log(1/\alpha_N)/m_N\to0$,
where $m_N=\Theta(N)$ is the per-bin sample size (no cap).  The null threshold
is $t_{\alpha_N}=\sqrt{2\kappa^2/m_N}(1+\sqrt{2\log(1/\alpha_N)})$, so
$\sqrt{m_N}\,t_{\alpha_N}\to\infty$ while
$\sqrt{m_N}\cdot O_P(B/\sqrt N)=O_P(1)$: same-class pairs have
$\widehat{\mmd}\le O_P(1/\sqrt{m_N})+O_P(B/\sqrt N)<t_{\alpha_N}$ eventually,
so they are merged, with probability $\to1$.  Cross-class pairs have a fixed
gap $\varepsilon>0$ while $t_{\alpha_N}\to0$, so by Fact~\ref{fact:dev} they
are split with probability $\to1$.  A union bound over the finitely many
pairs gives $\Pr[\hat\calQ_B=\Qstar_B]\to1$.

\paragraph{Under (a) alone.}
Population edges that are not in a gap fall inside some $I_i$; a bin that
contains a gap $G_i$ in its interior (no edge in $\overline{G_i}$) is a
\emph{straddling} bin.  There are $n-1$ gaps, so at most $n-1$ straddling
bins.  A straddling bin's law is $\sum_j\lambda_jP_j$ over the classes it
intersects with all $\lambda_j\in(0,1)$; for two consecutive classes,
$\mmd(\lambda P_i+(1-\lambda)P_{i+1},P_i)=(1-\lambda)\mmd(P_i,P_{i+1})>0$, so
it is distinct from every pure law and contributes at most one extra class.
If every class has mass at least $2/B$, every class also owns a pure bin, and
$n\le|\Qstar_B|\le2n-1$.  The recovery argument above applies verbatim to
$\Qstar_B$ (the straddling bins' mixture weights are $N$-independent).

\paragraph{Without (a): unbounded noise.}
If the class-conditional laws of $z$ have overlapping supports (e.g.\ Gaussian
noise), $F$ is strictly increasing, every bin contains observations of at least
two classes with fixed positive proportions, and the bin laws are proper
mixtures whose weights do not depend on $N$.  The bound
$\Pr[\text{cross}]\le2e^{-\Delta^2/(8\sigma^2)}$ for sub-Gaussian noise with
class-mean gap $\Delta$ bounds the \emph{proportion} of foreign observations
in a boundary bin by $\approx2Be^{-\Delta^2/(8\sigma^2)}$; this proportion is
constant in $N$, so a consistent test eventually splits the boundary bin (an
earlier version of this proposition claimed the proportion vanishes by
Borel--Cantelli, which is false).  Generically all bin laws are then distinct
and $|\Qstar_B|=\Khist$; the finite-sample estimate is a coarsening at the
test's resolution (Remark~4 of the main paper), which is the regime of all
Gaussian-noise experiments in the paper.

\paragraph{The environments of the paper.}
K-cycle: $z=o=s/K+\xi$, classes of equal mass $1/K$; with $B=K$ the population
edges are $F^{-1}(j/K)$, which by the symmetry of the equal-weight mixture lie
at the midpoints $(j-\tfrac12)/K$ up to exponentially small tail terms, so (b)
holds approximately while (a) fails ($\xi$ Gaussian).  Gaussian-mixture random
walk: uniform stationary law, $B=N_s$, same situation.  Tiger: symmetric
marginal, so the median edge is at $0$ and any even $B$ is aligned with the
sign boundary.  Chain and Dec-Tiger: $B=4$ aligned as above.

% ============================================================
\section{Experimental Details}
\label{sec:envs}
% ============================================================

Common defaults (\texttt{ContinuousEstimatorConfig}): history length
$\ell=1$, window $k=2$, $m_{\min}=20$, $m_{\max}=300$ continuation vectors per
side (deterministic per-pair subsampling), RBF kernel with the median heuristic
$2\sigma^2=\mathrm{median}\{\|x-y\|^2\}$ over at most 500 pooled points,
permutation $p$-value $(1+\#\{\ge\})/(R+1)$, Benjamini--Hochberg at level
$\alpha$ over all pairs.  Deviations are listed per experiment.

\subsection{Tiger (\texttt{envs/tiger\_continuous.py})}
States $\{\text{Left},\text{Right}\}$, uniform initial state; actions
$\{\text{listen},\text{open-left},\text{open-right}\}$; an open action re-draws
the state uniformly.  Observation $o=\mathrm{sgn}\cdot1+\mathcal N(0,\sigma^2)$
with $\mathrm{sgn}=-1$ (Left), $+1$ (Right), $\sigma=0.5$.  Uniform-random
policy, $T=20$.  Latent index $n=2$.  Estimator: $B=8$ (package default),
$\alpha=0.01$, $R=199$.  Main run: $N=3000$, 10 seeds
(\texttt{layer1\_ground\_truth/results/tiger\_continuous.csv}); $N$-sensitivity:
$N\in\{1000,3000,10000\}$, 5 seeds (\texttt{tiger\_N\_sensitivity.csv});
ablation: $N\in\{500,1000,2000,5000\}$, 3 seeds, $B=2$, $\alpha=0.05$, $R=99$
(\texttt{baselines/results/ablation\_sample\_efficiency.csv}).

\subsection{K-cycle (\texttt{envs/kcycle\_continuous.py})}
States $\{0,\dots,K-1\}$, deterministic increment mod $K$, uniform initial
state; observation $o=s/K+\mathcal N(0,0.1^2)$ (equally spaced means; the
$\sin(2\pi s/K)$ variant used in an earlier draft is not injective for even
$K$ and is not used); single null action; $T=50$.  Latent index $n=K$.
Estimator: $B=K$, $\alpha=0.01$, $R=199$, $N=3000$ ($5000$ for $K=16$),
10 seeds (\texttt{kcycle\_continuous.csv}).  Bin sensitivity: $B=2K$, 5 seeds
(\texttt{kcycle\_bins\_sensitivity.csv}).

\subsection{Gaussian-mixture random walk (\texttt{envs/gaussian\_mixture\_pomdp.py})}
$N_s\in\{3,4,5,6\}$ states, $s'=\mathrm{clip}(s+\delta,0,N_s-1)$,
$\delta\sim\mathrm{Unif}\{-1,0,1\}$; $o=s+\mathcal N(0,0.15^2)$; $T=50$,
$N=2000$; $\alpha=0.05$, $R=99$; $B=N_s$ (manual) and $B$ auto-selected by
silhouette over $\{2,\dots,32\}$ (Algorithm~2 below); 5 seeds
(\texttt{gaussian\_mixture.csv}).

\subsection{PO-HalfCheetah and PO-Ant (\texttt{layer2\_mujoco/})}
\texttt{HalfCheetah-v4}: 17-dim observation (8 joint positions, 9 joint
velocities); the wrapper zeroes indices $0$--$8$, leaving 8 velocity
coordinates.  \texttt{Ant-v4}: 27-dim observation; indices $13$--$26$ (14
velocities) are kept.  Zero-torque policy (and, for HalfCheetah, torques
$\propto\sin(\omega t+\phi_i)$ per joint); $N=5000$, $T=50$, $B=8$,
$\alpha=0.05$, $R=99$, 5 seeds (\texttt{halfcheetah\_continuous.csv},
\texttt{halfcheetah\_sinusoidal.csv}, \texttt{ant\_continuous.csv});
auto-bin: 3 seeds (\texttt{halfcheetah\_autobin.csv}).

Class statistics and value prediction are in
\path{halfcheetah_class_returns.csv},
\path{halfcheetah_value_prediction.csv} and
\path{halfcheetah_sinusoidal_value.csv}.  The binning scalar is
$z(o)=\|o\|_2$ (velocity norm).  No discretised G-test baseline is reported:
no result file exists for it.

\subsection{Dec-Tiger and chain Dec-POMDPs (\texttt{layer3\_rcrit/})}
\label{sec:rcrit-details}
Both agents run a listen-only policy so the latent state is constant within an
episode; $N=1000$, $T=20$, $B=4$, $m_{\min}=10$, $\alpha=0.05$, $R=99$.
Dec-Tiger: $o^i=\mathrm{sgn}+\mathcal N(0,\sigma_i^2)$; chain: 4-state random
walk (uniform stationary law), $o^i=s/4+\mathcal N(0,\sigma_i^2)$.  Each agent's
quotient is fitted independently; each step is labelled by the class of its
binned key; $\hat R_{\rm crit}=H(\hat Q_A,\hat Q_B)-H(\hat Q_B)$ (plug-in
entropies, clamped at $0$).

\paragraph{Analytical $H(Q_A\mid Q_B)$.}
Dec-Tiger: $Q_i=\mathrm{sign}(o^i)$ (the $B=4$ quantile edges of a symmetric
marginal put the middle edge at $0$), so $p_i=\Pr[Q_i\ne s]=\Phi(-1/\sigma_i)$.
With $\Pr[s=\pm1]=\tfrac12$,
\begin{align*}
  \Pr(Q_A=Q_B=s)     &=\tfrac12(1-p_A)(1-p_B)+\tfrac12p_Ap_B \quad(\text{each sign}),\\
  \Pr(Q_A\ne Q_B)    &=\tfrac12(1-p_A)p_B+\tfrac12p_A(1-p_B)\quad(\text{each pattern}),
\end{align*}
and $H(Q_A\mid Q_B)=H(Q_A,Q_B)-H(Q_B)$.  Symmetric $\sigma=0.5$:
$p=\Phi(-2)=0.0228$, $H(Q_A\mid Q_B)=0.2624$ bits.  Asymmetric
$(\sigma_A,\sigma_B)=(0.1,1.0)$: $p_A=\Phi(-10)\approx0$, $p_B=\Phi(-1)=0.1587$,
$H(Q_A\mid Q_B)=H_b(0.1587)=0.6311$ bits; ratio $2.41$.  Sweep values
($\sigma_A=0.1$): $\sigma_B=0.1,0.2,0.5,1.0,2.0\mapsto0.0000,0.0000,0.1566,0.6311,0.8915$.
(An earlier draft wrote these formulas with a threshold at $0.5$ and
$\Phi(-0.5/\sigma)$; that does not match the environment or the numbers.)
Chain: agent $i$ classifies $s$ by the nearest mean, i.e.\ intervals with edges
$(2s+1)/8$; $\Pr(Q_i=j\mid s)=\Phi((\text{hi}_j-s/4)/\sigma_i)-\Phi((\text{lo}_j-s/4)/\sigma_i)$;
joint over uniform $s$; values $\sigma_B=0.1,0.2,0.5,1.0\mapsto1.0584,1.4462,1.8576,1.9642$ bits.

\subsection{Baselines (\texttt{baselines/})}
G-test: quantile discretisation of the (scalar) observation into 3 or 8
symbols, then the finite-alphabet estimator of the companion paper
($\alpha=0.01$).  Wasserstein-$\varepsilon$: same continuation matrices,
$W_1$ between first-coordinate projections, merge iff $W_1<\varepsilon$.  DBC:
MLP encoder trained on $(d_{\rm embed}(h_1,h_2)-W_1(P_{h_1},P_{h_2}))^2$,
$k$-means with silhouette selection of $k$.  Result files:
\path{ablation_sample_efficiency.csv},
\path{bisim_comparison.csv} and \path{dbc_comparison.csv}.

% ============================================================
\section{Additional Results}
\label{sec:ablations}
% ============================================================

\subsection{Auto bin selection}

\begin{algorithm}[h]
\caption{\textsc{AutoBinSelection} (Algorithm~2)}
\label{alg:autobin}
\begin{algorithmic}[1]
\Require Trajectories, $B_{\min}=2$, $B_{\max}=32$, subsample size 2000.
\For{$B\in\{B_{\min},\dots,B_{\max}\}$}
  \State $k$-means with $B$ clusters on the subsampled binning scalars $z(o)$; silhouette score $s_B$.
\EndFor
\State \Return $\arg\max_Bs_B$.
\end{algorithmic}
\end{algorithm}

Silhouette measures cluster quality in \emph{observation} space and is a
heuristic proxy for the number of latent classes; it under-estimates $B$ when
classes share observation marginals.  Results: GaussMix---$B=N_s$ on all 20
runs, and the auto-$B$ estimate equals the manual one on every run
(\texttt{gaussian\_mixture.csv}); PO-HalfCheetah (passive)---$B=3$ on all 3
seeds with $\hat n=3,3,2$ (\texttt{halfcheetah\_autobin.csv}).

\subsection{K-cycle per-seed results and runtime}

\begin{table}[h]
\centering
\caption{K-cycle at $B=K$ (\texttt{kcycle\_continuous.csv}): all 60 runs return
  $\hat n=K$.  Wall-clock per run on one CPU core.}
\begin{tabular}{rrrrc}
\toprule
$K$ & Seeds & $\hat n$ & Runtime (s) & Correct \\
\midrule
2  & 0--9 & 2  & 2.9--3.6   & 10/10 \\
3  & 0--9 & 3  & 3.0--4.4   & 10/10 \\
4  & 0--9 & 4  & 4.0--5.1   & 10/10 \\
5  & 0--9 & 5  & 5.2--6.1   & 10/10 \\
8  & 0--9 & 8  & 8.0--11.7  & 10/10 \\
16 & 0--9 & 16 & 30.4--37.3 & 10/10 \\
\bottomrule
\end{tabular}
\end{table}

\begin{table}[h]
\centering
\caption{K-cycle at $B=2K$, 5 seeds (\texttt{kcycle\_bins\_sensitivity.csv}).
  Over-provisioned bins are split whenever the half-class bins' posteriors
  differ detectably ($3\le K\le8$); at $K=16$ the 32 bins are close enough that
  single-linkage chaining merges neighbours and the estimate falls below $K$.}
\label{tab:kcycle-2k}
\begin{tabular}{rrrl}
\toprule
$K$ & $B$ & $N$ & $\hat n$ per seed \\
\midrule
2 & 4 & 3000 & 2, 2, 2, 2, 2 \\
3 & 6 & 3000 & 3, 4, 3, 4, 4 \\
4 & 8 & 3000 & 8, 8, 8, 8, 8 \\
5 & 10 & 3000 & 10, 10, 10, 10, 10 \\
8 & 16 & 3000 & 16, 16, 16, 16, 16 \\
16 & 32 & 5000 & 12, 15, 10, 8, 15 \\
\bottomrule
\end{tabular}
\end{table}

\subsection{Tiger vs.\ $N$}
$B=8$, $\alpha=0.01$, $R=199$, 5 seeds (\texttt{tiger\_N\_sensitivity.csv}):
$\hat n=2$ on every seed at $N=1000$, $3000$ and $10000$.  Ablation runs
($B=2$, $\alpha=0.05$, $R=99$): $\hat n=2$ on 3/3 seeds at every
$N\in\{500,1000,2000,5000\}$.

\subsection{Dec-Tiger per-seed results}

\begin{table}[h]
\centering
\small
\caption{Dec-Tiger, $\alpha=0.05$, $R=99$, all seeds
  (\texttt{dec\_tiger\_asymmetric.csv}; seeds 0--4 of the symmetric block are
  also in \texttt{dec\_tiger\_rcrit.csv}).  Analytical: $0.2624$ (symmetric),
  $0.6311$ (asymmetric).}
\label{tab:rcrit-full}
\begin{tabular}{c|cccr|cccr}
\toprule
 & \multicolumn{4}{c|}{$(\sigma_A,\sigma_B)=(0.5,0.5)$} & \multicolumn{4}{c}{$(\sigma_A,\sigma_B)=(0.1,1.0)$}\\
Seed & $\hat n_A$ & $\hat n_B$ & $\hat R_{\rm crit}$ & rel.\ err. & $\hat n_A$ & $\hat n_B$ & $\hat R_{\rm crit}$ & rel.\ err. \\
\midrule
0 & 2 & 3 & 0.2662 & 1.5\%   & 2 & 4 & 0.5698 & 9.7\% \\
1 & 3 & 2 & 0.7651 & 191.6\% & 2 & 4 & 0.5567 & 11.8\% \\
2 & 2 & 2 & 0.2671 & 1.8\%   & 2 & 4 & 0.5544 & 12.2\% \\
3 & 2 & 3 & 0.2633 & 0.4\%   & 2 & 4 & 0.5545 & 12.1\% \\
4 & 2 & 2 & 0.2625 & 0.0\%   & 2 & 4 & 0.5654 & 10.4\% \\
5 & 3 & 3 & 0.7690 & 193.1\% & 3 & 4 & 1.0774 & 70.7\% \\
6 & 3 & 2 & 0.7607 & 189.9\% & 2 & 4 & 0.5700 & 9.7\% \\
7 & 3 & 3 & 0.7704 & 193.6\% & 2 & 4 & 0.5622 & 10.9\% \\
8 & 3 & 3 & 0.7624 & 190.5\% & 2 & 4 & 0.5489 & 13.0\% \\
9 & 2 & 2 & 0.2608 & 0.6\%   & 2 & 4 & 0.5530 & 12.4\% \\
\midrule
mean$\pm$std (all) & & & $0.515\pm0.264$ & 96.3\% & & & $0.611\pm0.164$ & 17.3\% \\
median (all)       & & & 0.514 & 95.8\%  & & & 0.559 & 12.0\% \\
mean$\pm$std ($\hat n_A=2$) & & & $0.264\pm0.003$ & 0.9\% & & & $0.559\pm0.008$ & 11.4\% \\
\bottomrule
\end{tabular}
\end{table}

Over-splitting of agent~A ($\hat n_A=3$) triples the estimate; over-splitting of
agent~B is harmless.  The mechanism is the resolution effect: with
$\sigma=0.5$ the posteriors of the two same-side bins differ by up to
$\approx0.05$ in posterior weight (the $B{=}4$ quantile bins of a $\pm1$, $\sigma{=}0.5$ marginal give $\Pr(\text{Right}\mid b)\approx0.0001$ and $0.045$ for the two left-side bins), at the edge of what $R=99$ permutations and
$m_{\max}=300$ resolve, and BH with only six $p$-values rejects all of them
once all six reach the floor $1/(R+1)=0.01$.  Table~\ref{tab:rcrit-alpha001}
reports the sensitivity run at $\alpha=0.01$ with $R=999$ (at $R=99$ the floor
equals $\alpha$ and every quotient collapses to one class).

\begin{table}[h]
\centering
\small
\caption{Dec-Tiger, $\alpha=0.01$, $R=999$, 10 seeds
  (\texttt{dec\_tiger\_sym\_alpha001\_R999.csv},
  \texttt{dec\_tiger\_asym\_alpha001\_R999.csv}).}
\label{tab:rcrit-alpha001}
\begin{tabular}{lcccc}
\toprule
Condition & $\hat n_A$ per seed & $\hat n_B$ per seed & $\hat R_{\rm crit}$ mean$\pm$std & median rel.\ err. \\
\midrule
$(0.5,0.5)$ & 2 (9 seeds), 3 (seed 6) & 2 (8), 3 (2) & $0.319\pm0.156$ (all), $0.270\pm0.007$ ($\hat n_A{=}2$) & 3.2\% \\
$(0.1,1.0)$ & 2 (10) & 4 (10) & $0.561\pm0.010$ & 11.4\% \\
\bottomrule
\end{tabular}
\end{table}

\subsection{Noise sweep per level}

\begin{table}[h]
\centering
\small
\caption{Sweep, $\sigma_A=0.1$, 5 seeds per level (\texttt{rcrit\_sweep.csv},
  \texttt{chain\_rcrit.csv}).  Relative errors are undefined when the
  analytical value is $0$; absolute errors are given.}
\begin{tabular}{lrrrrl}
\toprule
Env & $\sigma_B$ & analytical & $\hat R_{\rm crit}$ mean$\pm$std & mean abs.\ err.\ (bits) & $\hat n_A$ / $\hat n_B$ \\
\midrule
Dec-Tiger & 0.1 & 0.0000 & $0.081\pm0.055$ & 0.081 & 2 / 2 \\
Dec-Tiger & 0.2 & 0.0000 & $0.081\pm0.055$ & 0.081 & 2 / 2 \\
Dec-Tiger & 0.5 & 0.1566 & $0.180\pm0.016$ & 0.024 (median rel.\ 12.6\%) & 2 / \{2,3\} \\
Dec-Tiger & 1.0 & 0.6311 & $0.560\pm0.007$ & 0.071 (median rel.\ 11.8\%) & 2 / 4 \\
Dec-Tiger & 2.0 & 0.8915 & $0.878\pm0.021$ & 0.022 (median rel.\ 2.5\%) & 2 / \{2,3,4\} \\
\midrule
Chain & 0.1 & 1.0584 & $1.062\pm0.008$ & 0.007 (0.67\%) & 4 / 4 \\
Chain & 0.2 & 1.4462 & $1.450\pm0.007$ & 0.005 (0.33\%) & 4 / 4 \\
Chain & 0.5 & 1.8576 & $1.879\pm0.026$ & 0.026 (1.40\%) & 4 / \{2,3,4\} \\
Chain & 1.0 & 1.9642 & $2.000\pm0.000$ & 0.036 (1.82\%) & 4 / 1 \\
\bottomrule
\end{tabular}
\end{table}

\subsection{Neural encoder}
The code ships a GRU/InfoNCE history encoder
(\texttt{neural\_encoder.py}) trained on the MMD partition; no results for it
are reported in this paper because no result files were produced.

\subsection{PO-HalfCheetah class statistics and value prediction}
Tables~\ref{tab:halfcheetah_class_returns}--\ref{tab:sinval} follow.

\begin{table}[h]
\centering
\small
\begin{tabular}{crrrrr}
\toprule
Rank & Class & $\bar{v}_{\ell_2}$ & $\sigma_{v}$ & $\bar{G}$ & $\sigma_G$ \\
\midrule
  0 & 7 & 0.005 & 0.001 & -0.00 & 0.00 \\
  1 & 6 & 0.009 & 0.001 & 0.00 & 0.00 \\
  2 & 5 & 0.016 & 0.004 & 0.00 & 0.00 \\
  3 & 4 & 0.030 & 0.005 & 0.00 & 0.01 \\
  4 & 3 & 0.058 & 0.014 & 0.02 & 0.03 \\
  5 & 0 & 0.159 & 0.043 & -0.03 & 0.19 \\
  6 & 2 & 0.544 & 0.251 & 0.07 & 0.27 \\
  7 & 1 & 2.399 & 1.306 & -0.14 & 0.43 \\
\bottomrule
\end{tabular}
\caption{Per-class statistics for the $n{=}8$ quotient partition on
PO-HalfCheetah. Classes are sorted by mean velocity $\ell_2$ norm.
$\bar{G}$ = mean discounted return ($\gamma{=}0.99$) from each timestep
assigned to that class. Kruskal-Wallis $H=4897.3$, $p \approx 10^{-1055}$ (computed over time steps, which are dependent within trajectories; the $p$-value is not evidence of a trajectory-level effect).\vspace{-4pt}}
\label{tab:halfcheetah_class_returns}
\end{table}
\begin{table}[h]
\centering
\small
\begin{tabular}{lrrr}
\toprule
Model & Features & $R^2$ & MAE \\
\midrule
  Null (constant mean) & 0 & 0.000 & 0.087 \\
  Quotient classes ($n=3$, $\eta^2=0.010$) & 3 & 0.010 & 0.083 \\
  Velocity $\ell_2$ (oracle) & 1 & 0.066 & 0.084 \\
\bottomrule
\end{tabular}
\caption{Class-conditioned value prediction on PO-HalfCheetah (zero-action policy, $\gamma=0.99$, $n=51,000$ timesteps). The auto-selected $B=3$ quotient partition explains $1.0\%$ of return variance (linear $R^2$) and $\eta^2=0.010$ of total variance (ANOVA effect size). The Kruskal-Wallis $p$-value ($4.92\times10^{-117}$) treats time steps as independent and is not evidence of a trajectory-level effect.  All $R^2$ values here are in-sample; see Table~\ref{tab:hcsplit} for the held-out, trajectory-level comparison against a nonlinear speed baseline, which shows the classes carry no information beyond speed.}
\label{tab:halfcheetah_value_prediction}
\end{table}

\begin{table}[h]
\centering
\small
\begin{tabular}{lrr}
\toprule
Model & $R^2$ & MAE \\
\midrule
Null (constant mean) & 0.000 & 1.551 \\
Quotient classes (auto-selected $B$) & 0.163 & 1.426 \\
Velocity $\ell_2$ (oracle) & 0.082 & 1.492 \\
\bottomrule
\end{tabular}
\caption{Class-conditioned value prediction on PO-HalfCheetah under the
  sinusoidal locomotion policy ($\gamma=0.99$, $101{,}000$ time steps;
  \texttt{halfcheetah\_sinusoidal\_value.csv}).}
\label{tab:sinval}
\end{table}
\clearpage

% ============================================================
%  References
% ============================================================
\bibliographystyle{plainnat}
\bibliography{refs}

\end{document}
